%% ma: L11-B16-0061
%% lop: 11
%% chuong: 5
%% bai: 16
%% dang: 11.16.4
%% loai: DS
%% muc: [TH, TH, VD, VD]
%% dap_an: "ĐSĐS"
%% nguon: "tai-lieu-hoc-tap-mon-toan-lop-11-theo-chuong-trinh-giao-duc-pho-thong-2018.pdf (Lớp Toán Thầy Duy, Phạm Lê Duy, THPT Chu Văn An), Chương 3 Bài 2, Phần 2, Câu 5 (tr.378)"
%% kien_thuc: "Lớp 11 Bài 16 (giới hạn của hàm số)"
%% trang_thai: chinh-thuc
%% ghi_chu: "Đã duyệt 03/11/2026."
\begin{ex}
Cho các số thực $a,b,c$ thỏa mãn $c^2+a=18$ và $\lim\limits_{x\to+\infty}\left(\sqrt{ax^2+bx}-cx\right)=-2$. Khi đó
\choiceTF
{\True $a=9$.}
{$b=3c$.}
{\True $a=3c$.}
{$P=a+b+5c=14$.}
\loigiai{\sloppy Giới hạn hữu hạn nên $\sqrt a=c\ge0$, tức $a=c^2$; khi đó $\sqrt{ax^2+bx}-cx=\dfrac{bx}{\sqrt{ax^2+bx}+cx}\to\dfrac b{2c}=-2$, suy ra $b=-4c$. Từ $c^2+a=18$: $2c^2=18\Rightarrow c=3$, $a=9$, $b=-12$.
a) $a=9$. Đúng.
b) $b=-12\ne3c=9$. Sai.
c) $a=9=3c$. Đúng.
d) $P=9-12+15=12\ne14$. Sai.}
\end{ex}
